\begin{figure}[htbp]
\centering
\begin{tikzpicture}[font=\small, level distance=13mm,
  level 1/.style={sibling distance=36mm},
  level 2/.style={sibling distance=18mm},
  level 3/.style={sibling distance=10mm}]
  \tikzset{conf/.style={draw, rounded corners, fill=lightaccent, minimum width=18mm, minimum height=7mm, align=center},
           acc/.style={draw, rounded corners, fill=warnbg, minimum width=18mm, minimum height=7mm, align=center},
           edge from parent/.style={draw, -{Stealth[length=2mm]}}}

  \node[conf] {$c_0$}
    child { node[conf] {$c_{1,1}$}
      child { node[conf] {$c_{2,1}$} }
      child { node[conf] {$c_{2,2}$} }
    }
    child { node[conf] {$c_{1,2}$}
      child { node[conf] {$c_{2,3}$} }
      child { node[acc] {elfogadó\\ág} }
    }
    child { node[conf] {$c_{1,3}$}
      child { node[conf] {$c_{2,4}$} }
      child { node[conf] {$c_{2,5}$} }
    };

  \node[align=center, font=\scriptsize] at (0,-4.55) {A nemdeterminisztikus gép akkor fogad el, ha van legalább egy elfogadó számítási ág.};
\end{tikzpicture}
\caption{Nemdeterminisztikus Turing-gép számítási fája}
\label{fig:zv12_ndtm_szamitasi_fa}
\end{figure}
